\scanpage{175}

\uline{Construction}\quad Soit \deleted{\unsure{G}}
\[
(1)\qquad G=\mathrm{GL}(2,\widehat{\mathbb Z}).\widehat{\mathbb Z}
\qquad(\text{produit }\tfrac12\text{ direct})
\]
où $\mathrm{GL}(2,\widehat{\mathbb Z})$ opère sur $\widehat{\mathbb Z}$ par
\[
(2)\qquad u.\lambda=\det(u)\lambda
\]
\deleted{Considérons \unsure{l'hom.}}\inserted{Pour}
$u\in\mathrm{GL}(2,\widehat{\mathbb Z})$, on désigne par $u_0$ son image dans $G$
et pour $\lambda\in\widehat{\mathbb Z}$, on désigne par $\zeta(\lambda)$ son image dans $G$.

\[
(3)\qquad
\left\{
\begin{gathered}
\diagram{\begin{tikzcd}
\widetilde{\mathrm{SL}(2,\mathbb Z)}\arrow[r]&G\\
\widetilde\rho\arrow[r,mapsto]&\rho_0.\zeta(-2)\\
\widetilde\sigma\arrow[r,mapsto]&\sigma_0.\zeta(3)
\end{tikzcd}}\\
\left[
\diagram{\begin{tikzcd}
\widetilde\omega_0=\widetilde\rho^{-3}=\widetilde\sigma^2
 \arrow[r,mapsto]&\omega_0.\zeta(6)\\
\ell'_0=\widetilde\omega_0^2\arrow[r,mapsto]&\zeta(12)
\end{tikzcd}}
\right]\\
\diagram{\begin{tikzcd}
\widetilde{\unsure{\Gamma}}_0=\widetilde\tau'\arrow[r,mapsto]&
\deleted{\unsure{\zeta}}\unsure{\Gamma}_0
\end{tikzcd}}
\end{gathered}\right.
\]
\deleted{Soit} Le \unsure{sgroupe}
$\deleted{\unsure{\mathrm{SL}(2,\mathbb Z)}}\ \deleted{SG}$
\[
(4)\qquad SG=\mathrm{SL}(2,\widehat{\mathbb Z}).\widehat{\mathbb Z}
\]
image réciproque de $\mathrm{SL}(2,\widehat{\mathbb Z})$ dans $G$,
\deleted{opère} admet comme centre $\mu\times\widehat{\mathbb Z}$ -- il s'ensuit que l'opération de
$\deleted{\unsure{S}}G$ sur $SG$ par autom. int. se factorise en une opération de
$G/\widehat{\mathbb Z}^{\unsure{*}}\simeq\mathrm{GL}(2,\widehat{\mathbb Z})$
sur $SG$, donc si $\alpha\in\mathrm{GL}(2,\widehat{\mathbb Z})$, $x\in SG$,
$\alpha(x)=\alpha x\alpha^{-1}$ est défini \deleted{\unsure{par}} comme
$\widetilde\alpha x\widetilde\alpha^{-1}$, sans référence explicite à un relèvement parti-\newline
culier $\widetilde\alpha$ de $\alpha$ dans $G\deleted{\unsure{S}}$ :
\[
(5)\qquad
\alpha(x)\ \text{ou}\ \alpha x\alpha^{-1}
\overset{\mathrm{def}}=
\widetilde\alpha(x)\deleted{\unsure{\alpha}}
=\widetilde\alpha x\widetilde\alpha^{-1}
=\alpha_0(x)=\alpha_0x\alpha_0^{-1}.
\]
Considérons
\[
(6)\qquad
\left\{\begin{gathered}
\beta\in\mathrm{SL}(2,\widehat{\mathbb Z}).\deleted{\unsure{\mathrm{GL}}},
\\
\text{i.e.}\qquad
\beta=\underline u\deleted{\unsure{\beta}}(t\sigma+u)^{-1}
\end{gathered}\right.
\]
\margin{\[
\beta=\underline u b^{-1},\quad
\underline u=\begin{pmatrix}\mu&\lambda\\0&1\end{pmatrix},\quad
b\in\operatorname{Norm}_{\mathrm{GL}(2,\widehat{\mathbb Z})}(L_0)
\]
\deleted{\unsure{par un isom.}\ $\longrightarrow\ L_0$}
\[
\nu_2\in\{0,1\},\quad t^2+u^2=\mu\in\widehat{\mathbb Z}^{*}.
\qquad t,u\in\widehat{\mathbb Z}
\]}
Considérons un rel. $\widetilde\beta$ de $\beta$ dans $SG$
(p. ex. $\widetilde\beta_{\unsure{1}}=\beta_0$), d'où
\[
(7)\qquad
\widetilde\eta\overset{\mathrm{def}}=
\deleted{\unsure{\widetilde\beta(\widetilde\sigma)}}.
[\widetilde\beta,\widetilde\sigma]
=\widetilde\beta\widetilde\sigma(\widetilde\beta)^{-1}
=\widetilde\beta(\widetilde\sigma)
 \unsure{\widetilde\sigma^{-1}|\widetilde\beta^{-1}}
=\beta(\widetilde\sigma)\widetilde\sigma^{-1}.
\]
On a
\[
\sigma(\widetilde\eta)\overset{\unsure{(4)}}=
\widetilde\sigma(\widetilde\eta)
=\widetilde\sigma(\widetilde\beta)
 \underbrace{\widetilde\sigma^2(\widetilde\beta)^{-1}}=
\]
\margin{175}
